\section{Introduction}
Reinforcement learning from human feedback (RLHF) is a necessary but opaque tool underlying the success of popular language models (LMs) such as OpenAI's ChatGPT~\citep{ChatGPT} and Anthropic's Claude~\citep{bai2022training}. 
The prevalence of RLHF stems from its efficacy at circumventing one of the greatest difficulties in integrating human preferences into language models: specifying an explicit reward~\citep{christiano2017deep}. 
Reward models (RMs) are central to this process. 
They are created by copying the original language model and training it on labeled preference data, producing a model that can predict whether  one piece of text is likely to be preferred over another. 
A reinforcement learning optimizer then uses this reward model signal to update the parameters of the original model, improving performance on a variety of tasks~\citep{ouyang2022training, touvron2023llama}.

While the post-RLHF model (known as the \textit{policy}) and even the pretrained model are extensively documented and evaluated, the basic properties of the RLHF process like the RMs receive far less attention. 
% Though reward models are central to the effectiveness of RLHF and potentially glimpsing into how human values map to LMs, they remain under-evaluated. 
Recent work on training reward models~\citep{starling2023, llm-blender-2023} has begun to fill this gap, but utilizes validation sets from previous RLHF training processes, such as Anthropic's Helpful and Harmless data~\citep{bai2022training} or OpenAI's Learning to Summarize~\citep{stiennon2020learning}, which are known to have ceilings on accuracy between 60 and 70\% due to inter-annotator disagreement~\citep{wang2024secrets}. 
% Similar investigations have yet to be conducted for Direct Policy Optimization (DPO) models. 
Moreover, newly released preference data aiming to expand the diversity of preference training datasets such as UltraFeedback~\citep{cui2023ultrafeedback}, UltraInteract~\citep{yuan2024advancing} and Nectar~\citep{starling2023}, do not have test sets, necessitating a new style of evaluation for RMs.  

We begin to rectify the lack of evaluation methods by introducing~\name, the first toolkit for benchmarking reward models. 
RLHF is a broadly applicable process used to enhance specific capabilities of LMs such as safety~\citep{dai2023safe} or reasoning~\citep{lightman2023let, havrilla2024teaching} as well as general capabilities such as instruction following~\citep{ouyang2022training} or ``steerability''~\citep{askell2021general,bai2022training}. 
Thorough evaluations of RMs will also cover these categories. 
In this work, we curate data to create structured comparisons across a variety of reward model properties.
Each sample is formatted as a prompt with a human-verified chosen and rejected completion.
We design subsets so as to vary in difficulty and coverage. 
Some subsets are solved by small RMs, reaching 100\% accuracy, but others are harder to differentiate and still have state-of-the-art performance around 75\%, with many models around the random baseline.
% Some are constructed such that many reward models can differentiate chosen from rejected completions, reaching nearly 100\% accuracy. Others are more difficult and state-of-the-art language models only reach the 60 to 70\% range.

We aim to map the current landscape of openly available reward models via a leaderboard for~\name.
We have evaluated over 80 models, such those trained as classifiers, including UltraRM~\citep{cui2023ultrafeedback}, Starling~\citep{starling2023}, PairRM~\citep{llm-blender-2023}, SteamSHP~\citep{pmlr-v162-ethayarajh22a}, models from Reward rAnked FineTuning (RAFT)~\citep{dong2023raft}, and others. 
We also evaluate popular chat models trained with Direct Policy Optimization (DPO)~\citep{rafailov2023direct}, for example, Zephyr-$\beta$~\citep{tunstall2023zephyr}, Qwen-Chat~\citep{bai2023qwen}, StableLM~\citep{bellagente2024stable}, and T\"ulu 2~\citep{ivison2023camels} to ground recent debates on RLHF methods and showcase specific datasets where they fall short.

With these models, we compare scaling, test reasoning capabilities, highlight three buckets of refusal behavior, and share more details on the inner workings of RMs.
The accompanying code-base %\footnote{Code available here: \url{https://github.com/allenai/reward-bench}.} 
provides a common inference stack for many variations of models and we release many text-score pairs to analyze their performance.
With \textbf{\name}, we:
\begin{enumerate}[leftmargin=0.5cm]
    \item Release a common \textbf{framework for evaluating the many different architectures of reward models}, along with tools for visualization, training, and other analysis.
    We also release all data used in the evaluation, composed of text-score pairs for all inputs, to enable further data analysis on the properties of reward models.\footnote{Data is here: \url{https://huggingface.co/datasets/allenai/reward-bench-results}.}
    \item Illustrate the \textbf{differences between DPO and classifier-based reward models} across a variety of datasets. 
    DPO models, while more plentiful due to the method's simplicity, fail to generalize to popular preference data test sets and present a higher variance in performance.
    \item Chart the \textbf{landscape of current state-of-the-art reward models}. 
    We showcase the scaling laws, the propensity to refuse (or not), the reasoning capabilities, and more for popular RMs.
    \item Show the \textbf{limitations of existing preference data test sets} for evaluating these models, showcasing common pitfalls of RMs on subtle, but challenging instruction pairs (e.g. intentionally modified rejected responses, which superficially look high quality but answer the wrong prompt).
\end{enumerate}

% We hope this benchmark enables more advanced reward model training, scientific understanding of the integration of human preferences in LMs, and ultimately better aligned, open language models.
